\documentclass[12pt, a4paper, oneside]{article}
\usepackage{amsmath, amsthm, amssymb, bm, graphicx, hyperref, mathrsfs, booktabs}
\usepackage[margin=2.5cm]{geometry}

\linespread{1.3}
\newtheorem{theorem}{Theorem}
\newtheorem{definition}[theorem]{Definition}
\newtheorem{remark}[theorem]{Remark}

\newcommand{\rvec}{\mathbf{r}}
\newcommand{\mvec}{\mathbf{m}}
\newcommand{\avec}{\mathbf{a}}
\newcommand{\Na}{N}
\newcommand{\Ntri}{N_\triangle}
\newcommand{\eps}{\varepsilon}

\begin{document}

\thispagestyle{empty}
\vspace*{\fill}
\begin{center}
    \Huge\textbf{Topological Charge Density\\on the Lattice}

    \vspace{1cm}
    \Large\textbf{Algorithm and Formulas}

    \vspace{0.8cm}
    \large \texttt{corr/tools/topocharge.py}
\end{center}
\vspace*{\fill}

\newpage
\tableofcontents
\newpage

% ======================================================================
\section{Overview}
% ======================================================================

The topological charge of a two-dimensional spin texture is

\begin{equation}
    Q = \frac{1}{4\pi}\int \mvec\cdot
        \left(\frac{\partial\mvec}{\partial x}\times
              \frac{\partial\mvec}{\partial y}\right)\,dx\,dy ,
    \label{eq:continuum}
\end{equation}

with $\mvec$ the normalised magnetisation. It counts how many times the texture
wraps the unit sphere, and is therefore an integer for any configuration on a
closed surface.

A lattice has no derivatives. Discretising Eq.~\eqref{eq:continuum} by finite
differences would produce a quantity that is only approximately integer and
that drifts with thermal noise --- the integrality, the one property worth
having, would be the first casualty. The standard construction of Berg and
L\"uscher instead covers the plane with elementary, non-overlapping,
anticlockwise-ordered triangles and assigns to each the signed area $A_l$ that
its three spins span on the unit sphere:

\begin{equation}
\boxed{
    Q = \frac{1}{4\pi}\sum_l A_l .
}
\label{eq:berg}
\end{equation}

This is exact rather than approximate: it is the area of the image of a
piecewise-geodesic interpolation of the spin field, so $Q$ is the degree of a
map from a closed surface to $S^2$ and comes out an exact integer
(Sec.~\ref{sec:integrality}).

What this module returns is not $Q$ but the \emph{density} behind it: each
triangle's $A_l/4\pi$ shared equally among its three vertices, giving one
number $q_i$ per site with

\begin{equation}
    \sum_i q_i = Q .
    \label{eq:sumrule}
\end{equation}

The density is the object of interest --- it shows where the charge sits, which
is what a picture of a skyrmion is --- and Eq.~\eqref{eq:sumrule} then serves
as a free and extremely sharp correctness check on the whole path.

\begin{remark}[Code layout]
\texttt{tools/topocharge.py} holds the numerics and nothing else: no argument
parsing, no plotting. The triangulation is built by \texttt{build\_triangles},
the areas by \texttt{triangle\_charge}, and \texttt{site\_density} is the only
function the renderer calls. The result is exposed as one more entry in the
\texttt{DERIVED} table of \texttt{tools/frame.py}, so it reaches
\texttt{tools/animate.py} without a line of new rendering code. The self-test
is in the module's \texttt{\_\_main\_\_} block; run the file directly.
\end{remark}

% ======================================================================
\section{Input}
% ======================================================================

One frame, as returned by \texttt{tools/dumpframe.py}: the site positions
$\rvec_i$, a three-component vector per site read from named dump columns, and
the box bounds. The box is the reason \texttt{dumpframe} carries the
\texttt{BOX BOUNDS} block rather than discarding it --- the triangulation has to
wrap, and the periodic images are not otherwise known. It is converted to
lattice vectors by \texttt{geometry.lammps\_box\_to\_lattice}, shared with the
$\mathbf{q}$-resolved route, so a triclinic cell is handled the same way in
both.

\subsection{Two Requirements on the Sites}

\textbf{Only the magnetic sublattice.} Non-magnetic species carry no spin
direction and are not part of the texture. In a NiI$_2$ dump, two thirds of the
atoms are iodine; leaving them in makes the triangulation meaningless. Select
with \texttt{-\/-element}, or drop zero vectors with
\texttt{-\/-drop-zero-vector}.

\textbf{One site per primitive cell, on a regular $n_1\times n_2$ grid.} This
is what \texttt{replicate $n_1$ $n_2$ 1} of a cell with a single magnetic atom
produces, and it is what makes the elementary triangles of Sec.~\ref{sec:tri}
well defined. The grid is read off the positions (Sec.~\ref{sec:grid}) rather
than assumed, and a
mismatch between $n_1n_2$ and the number of sites is reported as an error
naming the most likely cause, rather than silently triangulating something
else.

\subsection{Normalisation}

Eq.~\eqref{eq:continuum} is written for a unit vector field, and the identities
of Sec.~\ref{sec:area} assume $\lVert\mvec_i\rVert=1$. The spins are therefore
normalised on input,

\begin{equation}
    \mvec_i = \frac{\mathbf{S}_i}{\lVert\mathbf{S}_i\rVert} ,
\end{equation}

and a site with $\mathbf{S}_i=0$ --- which has no direction at all --- raises an
error instead of producing a silent \texttt{NaN}.

% ======================================================================
\section{Triangulation}
\label{sec:tri}
% ======================================================================

\subsection{Why Not Delaunay}

The obvious approach --- Delaunay-triangulate the positions each frame --- is
wrong here for a specific reason. On a perfect triangular lattice every
elementary rhombus has its four corners on a common circle, so the Delaunay
triangulation is degenerate: either diagonal is admissible, and which one a
library returns is decided by rounding. Thermal displacement then makes the
choice flip from frame to frame. The charge itself survives this (both
triangulations are valid coverings), but the \emph{density} does not: $q_i$
depends on which triangles a site belongs to, so the picture flickers. A
per-frame Delaunay is also $O(\Na\log\Na)$ against $O(\Na)$ for the
construction below.

The connectivity is instead built once and reused for every frame. This
requires the atom order to be identical in every frame, which LAMMPS guarantees
under \texttt{dump\_modify $\langle$id$\rangle$ sort id}.

\subsection{Reading the Grid off the Positions}
\label{sec:grid}

Let $f^{(1)}_i, f^{(2)}_i$ be the fractional coordinates of site $i$,

\begin{equation}
    \mathbf{f}_i = \rvec_i\,\mathbf{L}^{-1},
\end{equation}

with $\mathbf{L}$ the cell matrix whose rows are $\avec_1,\avec_2,\avec_3$.
Working in fractional coordinates rather than in $x,y$ is what makes a triclinic
cell --- the NiI$_2$ supercell has $xy = -99.48$~\AA{} --- no different from an
orthogonal one.

On a perfect $n_1\times n_2$ grid the first fractional coordinate takes the
values $f^{(1)} = k/n_1$. Hence the order parameter

\begin{equation}
\boxed{
    \Phi(n) = \left|\frac{1}{\Na}\sum_i e^{2\pi i\,n f^{(1)}_i}\right|
}
\label{eq:sharp}
\end{equation}

satisfies $\Phi(n_1) = 1$, and $\Phi(n) = 0$ for every $n < n_1$, since
$\sum_k e^{2\pi i nk/n_1}$ vanishes unless $n_1 \mid n$. The smallest $n$ with
$\Phi(n) > 0.9$ is therefore $n_1$; the same applied to $f^{(2)}$ gives $n_2$.

This is preferred to clustering the coordinates with a distance tolerance
because it degrades gracefully and needs nothing tuned. Thermal displacement of
r.m.s.\ fractional size $\sigma$ only lowers the peak,

\begin{equation}
    \Phi(n_1) \simeq \exp\!\left[-\tfrac{1}{2}(2\pi n_1\sigma)^2\right] ,
\end{equation}

so the criterion holds as long as the displacement stays well below the lattice
spacing --- exactly the regime in which "one site per cell" is meaningful at
all. Outside it the routine raises rather than guesses, and
\texttt{-\/-topo-grid $n_1$ $n_2$} overrides.

\subsection{Removing the Grid Offset}

The cell index follows from $i \approx n_1 f^{(1)}$, but rounding that directly
is fragile: the dump's origin is arbitrary, and if the grid happens to sit
half-way between integers every site rounds ambiguously. The offset is
therefore measured from the phase of the same sum and removed first,

\begin{equation}
\boxed{
    i = \left[\operatorname{round}\!\left(n_1 f^{(1)}
        - \frac{1}{2\pi}\arg\sum_j e^{2\pi i\,n_1 f^{(1)}_j}\right)\right]
        \bmod n_1 ,
}
\end{equation}

which centres every cluster of sites on an integer before rounding. The result
cannot depend on where the box origin sits. A collision --- two sites in one
cell --- is detected and reported.

\subsection{Elementary Triangles}

Each cell $(i,j)$ contributes the two triangles that tile its rhombus:

\begin{equation}
    \bigl[(i,j),\ (i{+}1,j),\ (i,j{+}1)\bigr],
    \qquad
    \bigl[(i{+}1,j),\ (i{+}1,j{+}1),\ (i,j{+}1)\bigr],
\end{equation}

with the indices taken modulo $n_1$ and $n_2$. There are exactly
$\Ntri = 2n_1n_2 = 2\Na$ triangles; they are non-overlapping, they cover the
cell once, and because the indices wrap the covered surface is a closed torus.
That closure is what Sec.~\ref{sec:integrality} rests on. In the code the whole
construction is four \texttt{np.roll} calls on the $(n_1,n_2)$ table of site
indices.

\subsection{Orientation}

Eq.~\eqref{eq:berg} requires the triangles to be ordered anticlockwise; ordering
one clockwise flips the sign of its $A_l$. Rather than test each triangle, note
that the vertex order above is anticlockwise precisely when the cell itself is,
so a single test on the cell vectors settles all $\Ntri$ at once:

\begin{equation}
    (\avec_1\times\avec_2)\cdot\hat z < 0
    \;\Longrightarrow\;
    \text{reverse every triangle.}
\end{equation}

\subsection{Layers}

A bilayer has two decoupled textures, and stitching them into one surface would
connect sites that are not neighbours. The sites are therefore split at the
midpoint of the $z$ range --- the same rule \texttt{frame.py} uses to draw its
two panels --- and each layer is triangulated and closed independently, so
Eq.~\eqref{eq:sumrule} holds per layer and the total is the sum of two
integers. \texttt{-\/-single-layer} treats the frame as one surface.

% ======================================================================
\section{The Signed Spherical Area}
\label{sec:area}
% ======================================================================

\subsection{The Two Identities}

For a triangle with unit spins $\mvec_i,\mvec_j,\mvec_k$, write

\begin{equation}
    a = \mvec_i\cdot\mvec_j, \qquad
    b = \mvec_j\cdot\mvec_k, \qquad
    c = \mvec_k\cdot\mvec_i, \qquad
    \rho = \sqrt{2(1+a)(1+b)(1+c)} .
\end{equation}

The signed area $A_l$ of the spherical triangle they span is fixed by the pair

\begin{equation}
\boxed{
    \cos\frac{A_l}{2} = \frac{1 + a + b + c}{\rho},
    \qquad
    \sin\frac{A_l}{2} = \frac{\mvec_i\cdot(\mvec_j\times\mvec_k)}{\rho} .
}
\label{eq:pair}
\end{equation}

The second identity is the origin of the sign rule usually quoted alongside the
first,
$\operatorname{sign}(A_l) = \operatorname{sign}[\mvec_i\cdot(\mvec_j\times\mvec_k)]$:
since $\rho>0$, the sign of $A_l$ is the sign of the scalar triple product,
which is in turn the orientation of the three spins on the sphere.

\subsection{Evaluation}

The two identities of Eq.~\eqref{eq:pair} share the denominator $\rho$, so
dividing them removes it:

\begin{equation}
\boxed{
    A_l = 2\operatorname{atan2}
    \bigl(\,\mvec_i\cdot(\mvec_j\times\mvec_k),\;\;
             1 + a + b + c\,\bigr) \in (-2\pi, 2\pi] .
}
\label{eq:atan2}
\end{equation}

This is what the code evaluates. It is three vectorised lines over all
$\Ntri$ triangles at once, it needs no square root, and the two-argument
arctangent recovers both the magnitude and the sign in a single call --- the
sign rule is not applied separately because it is already contained in the sign
of the first argument.

\subsection{Conditioning: Why Not \texorpdfstring{$\arccos$}{arccos}}
\label{sec:conditioning}

Taking $\arccos$ of the first identity alone and multiplying by the sign rule is
algebraically identical and numerically worse. Two effects, both verified
numerically in Sec.~\ref{sec:validation}:

\textbf{Amplification at small $A_l$.} From
$\frac{d}{dA}\cos\frac{A}{2} = -\frac{1}{2}\sin\frac{A}{2} \simeq -A/4$, an
absolute error $\delta$ in the argument becomes

\begin{equation}
    \delta A_l \simeq \frac{4\delta}{A_l} ,
\end{equation}

which diverges as $A_l\to0$. A smooth texture puts nearly every triangle there:
$A_l \sim (a/R)^2$ for a texture of length scale $R$ on a lattice of spacing
$a$.

\textbf{An absolute error floor.} The divergence is capped only by the
resolution of the argument itself. The largest double below $1$ differs from it
by $\delta = \eps/2 = 1.11\times10^{-16}$, and

\begin{equation}
    \arccos(1-\delta) \simeq \sqrt{2\delta} ,
\end{equation}

so a triangle whose true $A_l$ is zero returns

\begin{equation}
\boxed{
    |\delta A_l| \;\gtrsim\; 2\sqrt{2\delta} \;=\; 3.0\times10^{-8}\ \text{rad},
}
\label{eq:floor}
\end{equation}

no matter how smooth the texture is. The amplification from argument to result
is $1/\sqrt{2\delta}\approx 7\times10^{7}$. Eq.~\eqref{eq:atan2} has no such
floor: with $\sin$ and $\cos$ supplied separately, $\operatorname{atan2}$ is
well conditioned through $A_l = 0$.

\begin{remark}[How much this actually matters]
Eq.~\eqref{eq:floor} is a floor of $3.0\times10^{-8}$ rad per triangle, or
$2.4\times10^{-9}$ per site in $q_i$. A skyrmion peaks near
$q_i \sim 5\times10^{-3}$, so the $\arccos$ form would still produce a
recognisable picture --- the floor sits some six orders below the signal. The
argument for Eq.~\eqref{eq:atan2} is not that the alternative fails, but that it
is simpler (no $\rho$, no square root, no separate sign) \emph{and} exact, so
there is nothing to trade. The $\arccos$ form is retained in the self-test as an
independent cross-check of Eq.~\eqref{eq:atan2}.
\end{remark}

\subsection{Degenerate Triangles}

If any two spins of a triangle are exactly antiparallel, say $\mvec_j=-\mvec_i$,
then $1+a = 0$ and $\rho = 0$. Both sides of Eq.~\eqref{eq:pair} become $0/0$:
the numerator of the first is $b + c = \mvec_k\cdot(\mvec_i+\mvec_j) = 0$ and
that of the second vanishes with it. This is not a numerical accident but a
genuine degeneracy --- three spins with two antipodal do not determine a
spherical triangle, and $A_l$ is undefined.

The $\arccos$ form returns \texttt{NaN}. Eq.~\eqref{eq:atan2} returns
$\operatorname{atan2}(0,0) = 0$, which is a definite value for an undefined
quantity, so neither form is safe in this limit and the configuration itself is
what has to be checked. The self-test reports
$\min_{\langle ij\rangle}(1 + \mvec_i\cdot\mvec_j)$ so proximity to the
degeneracy is visible. Smoothly modulated textures --- spirals, skyrmions,
spin waves --- have small angles between neighbours and stay far from it; a
sharp antiferromagnet on the same triangulation would not, and would need the
magnetic unit cell treated as the repeating unit instead.

% ======================================================================
\section{From Charge to Density}
% ======================================================================

Eq.~\eqref{eq:berg} lives on triangles, but \texttt{tripcolor} with Gouraud
shading --- the renderer \texttt{frame.py} already uses --- wants one value per
\emph{vertex}. Each triangle's charge is therefore split equally among its
three sites:

\begin{equation}
\boxed{
    q_i = \frac{1}{4\pi}\sum_{l \ni i} \frac{A_l}{3} ,
    \qquad \sum_i q_i = \frac{1}{4\pi}\sum_l A_l = Q .
}
\label{eq:density}
\end{equation}

The second equality is immediate --- every triangle has exactly three vertices,
so the thirds reassemble --- and it is the sum rule of Eq.~\eqref{eq:sumrule}.
On the uniform triangulation each site belongs to six triangles.

\begin{remark}[Units]
$q_i$ is dimensionless and sums to the integer $Q$. The integrand of
Eq.~\eqref{eq:continuum} is a density per unit area; it is recovered as
$q_i/A_{\text{site}}$ with
$A_{\text{site}} = |\avec_1\times\avec_2|/\Na$ the area per site. Since that is
a global constant the two differ only by the scale on the colour bar, and the
dimensionless form is used because it makes the sum rule readable directly off
the data.
\end{remark}

\begin{remark}[The alternative]
Plotting $A_l/4\pi$ per triangle with \texttt{shading="flat"} is closer to
Eq.~\eqref{eq:berg} as written, but produces $2\Na$ facets instead of $\Na$
values and cannot be registered against the arrows drawn at the sites. The
vertex form is used for both reasons.
\end{remark}

% ======================================================================
\section{Integrality}
\label{sec:integrality}
% ======================================================================

\begin{theorem}
For any configuration of unit spins on the closed triangulation of
Sec.~\ref{sec:tri}, with no degenerate triangle, $Q$ as defined by
Eq.~\eqref{eq:berg} is an integer.
\end{theorem}

Interpolate the spins geodesically across each triangle. This maps the closed
surface covered by the triangulation --- a torus, because the cell indices wrap
--- continuously onto $S^2$. The total signed area swept is $4\pi$ times the
number of times the sphere is covered, i.e.\ $4\pi$ times the degree of the map,
and the degree of a map between closed oriented surfaces is an integer. No
smoothness of the texture is needed: the statement holds for a random
configuration just as much as for a skyrmion.

This is the single most useful property of the construction, because it fails
loudly if anything upstream is wrong. Every one of the following breaks it:

\begin{itemize}
    \item a triangle ordered clockwise while its neighbours are anticlockwise;
    \item indices that do not wrap, leaving the surface open at the boundary;
    \item a missing or duplicated triangle;
    \item two layers stitched into one surface;
    \item non-magnetic sites left among the magnetic ones.
\end{itemize}

Checking $\sum_i q_i \in \mathbb{Z}$ therefore tests the triangulation, the
orientation, the wrapping, the layer split and the site selection at once, on
whatever data is at hand --- no reference result required. It is the first thing
the self-test does, on random spins, where nothing else could serve as a
reference.

% ======================================================================
\section{Validation}
\label{sec:validation}
% ======================================================================

The self-test (\texttt{python3 tools/topocharge.py}) runs on ideal triangular
lattices with a triclinic cell, at two grid shapes to catch anything that
happens to work only when $n_1 = n_2$.

\begin{center}
\begin{tabular}{lll}
    \toprule
    Check & Result & \\
    \midrule
    Triangle count & $2n_1n_2$ & the rhombus tiling is complete \\
    Random spins, $24\times24$ & $Q = -3.000000000000$ & integrality, no reference needed \\
    Random spins, $30\times18$ & $Q = -6.000000000000$ & \\
    Skyrmion, both shapes & $Q = -1.000000000000$ & exactly one quantum \\
    Bilayer, one each & $Q = -2.000000000000$ & layers not stitched \\
    Jittered positions & identical connectivity & Eq.~\eqref{eq:sharp} survives noise \\
    $\arccos$ vs.\ Eq.~\eqref{eq:atan2} & $1.6\times10^{-14}$, $2.4\times10^{-9}$ & see below \\
    $\min(1+\mvec_i\cdot\mvec_j)$ & $1.87$, $1.91$ & far from degeneracy \\
    \bottomrule
\end{tabular}
\end{center}

The two figures in the $\arccos$ row are Sec.~\ref{sec:conditioning}'s two
regimes, and were traced to the individual triangles responsible. At
$24\times24$ the worst triangle has $A_l = 2.87\times10^{-3}$, and
$4\delta/A_l = 3.1\times10^{-13}$ against an observed $2.1\times10^{-13}$ --- the
amplification law. At $30\times18$ the worst has $A_l = -1.5\times10^{-17}$,
effectively zero, with $\cos(A_l/2)$ one ulp below $1$; Eq.~\eqref{eq:floor}
predicts $2.98\times10^{-8}$ rad and the observed discrepancy is
$2.98\times10^{-8}$ rad. Both mechanisms are quantitative, not qualitative.

\subsection{End-to-End}

Two further checks run through the full path, dump file included:

\textbf{A synthetic skyrmion dump} (40$\times$40, triclinic, two frames, the
second with the skyrmion displaced) read by \texttt{dumpframe} and reduced by
\texttt{site\_density} gives $Q = -1.000000000000$ on both frames. This
exercises the \texttt{BOX BOUNDS} parsing and the reuse of the cached
connectivity on a texture that has moved.

\textbf{\texttt{examples/s4.lammpstrj}} gives $\max_i|q_i| = 0$ exactly. That
trajectory is a single-$\mathbf{q}$ spin wave, modulated along one direction
only, so $\partial_x\mvec \parallel \partial_y\mvec$ and the integrand of
Eq.~\eqref{eq:continuum} vanishes identically. A one-dimensional texture
carrying no topological density is a statement about the physics, not about the
lattice, and the code reproduces it to the last bit.

% ======================================================================
\section{Usage}
% ======================================================================

The density is exposed as the derived colour \texttt{topo}, so it is drawn by
the existing renderer:

\begin{verbatim}
python3 tools/frame.py dump.lammpstrj --frame 1500 \
    --vector 'c_outsp[1]' 'c_outsp[2]' 'c_outsp[3]' \
    --color topo --element Ni --single-layer \
    --cmap RdBu_r --vmin -0.006 --vmax 0.006 --out topo.png
\end{verbatim}

\texttt{tools/animate.py} takes the same options and produces the animation.
Four practical points:

\begin{itemize}
    \item \texttt{-\/-vector} is required. \texttt{topo} is a function of the
    spin direction, so the three spin columns must be named even when no arrows
    are drawn (\texttt{-\/-no-arrows}).
    \item \texttt{-\/-element} (or \texttt{-\/-drop-zero-vector}) is required
    whenever the dump holds more than the magnetic sublattice.
    \item Set \texttt{-\/-vmin}/\texttt{-\/-vmax} by hand, symmetrically, with a
    diverging colormap. The density is strongly peaked at a skyrmion core, so
    an automatic range is set by that peak and leaves the surroundings flat.
    \item \texttt{-\/-topo-grid $n_1$ $n_2$} is only needed if
    Eq.~\eqref{eq:sharp} fails to identify the grid, which it reports rather
    than guessing.
\end{itemize}

% ======================================================================
\section{Summary of the Algorithm}
% ======================================================================

Steps 1--5 run once; steps 6--8 run per frame.

\begin{enumerate}
    \item \textbf{Convert} the box bounds to lattice vectors and the positions
    to fractional coordinates.
    \item \textbf{Split} into layers at the midpoint of the $z$ range, unless
    \texttt{-\/-single-layer}.
    \item \textbf{Identify} the grid $n_1\times n_2$ from Eq.~\eqref{eq:sharp};
    verify $n_1n_2 = \Na$.
    \item \textbf{Index} each site by cell, with the grid offset removed;
    verify no collisions.
    \item \textbf{Build} $2\Na$ triangles by wrapping the cell indices, and
    orient them from $(\avec_1\times\avec_2)\cdot\hat z$. Cache.
    \item \textbf{Normalise} the spins to unit length.
    \item \textbf{Evaluate} $A_l$ for every triangle by Eq.~\eqref{eq:atan2}.
    \item \textbf{Accumulate} $A_l/12\pi$ onto each of the triangle's three
    vertices, Eq.~\eqref{eq:density}.
\end{enumerate}

The per-frame cost is $O(\Na)$ with no transcendental work beyond one
$\operatorname{atan2}$ per triangle.

% ======================================================================
\section{Notation}
% ======================================================================

\begin{table}[h]
    \centering
    \begin{tabular}{ll}
        \toprule
        Symbol & Meaning \\
        \midrule
        $\Na$ & Magnetic sites in the layer, $= n_1 n_2$ \\
        $n_1, n_2$ & Cell repeats along $\avec_1, \avec_2$ \\
        $\Ntri$ & Triangles, $= 2\Na$ per layer \\
        $\mathbf{L}$ & Cell matrix, rows $\avec_1,\avec_2,\avec_3$ \\
        $\mathbf{f}_i$ & Fractional coordinates of site $i$ \\
        $\mvec_i$ & Unit spin on site $i$ \\
        $a,b,c$ & Pairwise dot products within a triangle \\
        $\rho$ & $\sqrt{2(1+a)(1+b)(1+c)}$, the shared denominator \\
        $A_l$ & Signed spherical area of triangle $l$, Eq.~\eqref{eq:atan2} \\
        $q_i$ & Topological charge density on site $i$, Eq.~\eqref{eq:density} \\
        $Q$ & Total topological charge, $\sum_i q_i \in \mathbb{Z}$ \\
        $\eps$ & Machine epsilon, $2.22\times10^{-16}$ \\
        \bottomrule
    \end{tabular}
    \caption{Notation.}
\end{table}

\end{document}
